% AUTO-GENERATED from fhops_operational_formulation.md -- do not edit directly.
FHOPS' deterministic operational solver is formulated on a day-shift
grid and maximizes weighted delivered production while penalizing
leftovers, landing over-capacity slack, and machine movement costs. The
equations below state the FHOPS 1.0.1 model and mirror the implemented
Pyomo model in
\texttt{fhops.model.milp.operational.build\_operational\_model} and the
bundle normalization in
\texttt{fhops.model.milp.data.build\_operational\_bundle}.

\textbf{Problem statement.} Given harvest blocks, machine roles, shift
calendars, block windows, landing capacities, and harvest-system role
prerequisites, choose machine-block assignments and per-shift production
quantities to maximize weighted production subject to feasibility and
sequencing constraints.

\textbf{Labels.} For traceability, the objective is tagged \textbf{OBJ},
the optional earliness stage \textbf{OBJ2}, the constraint blocks
\textbf{E1}--\textbf{E13}, the optional initial state \textbf{INIT}, and
the domain declarations \textbf{D1}. Blocks that persist from FHOPS
1.0.0 keep their 1.0.0 labels (\textbf{E1}--\textbf{E11}). Of these,
\textbf{E6} (machine moves), \textbf{E7} (staged inventory), \textbf{E8}
(activation and head start, now \textbf{E8a}--\textbf{E8c}), \textbf{E9}
(loader truckload threshold, which replaces the 1.0.0 batching
variables) and \textbf{E11} (landing capacity per shift slot) were
reformulated in 1.0.1. \textbf{E12} (remaining output), \textbf{E13}
(locked assignments), \textbf{INIT}, and \textbf{OBJ2} are new.

\textbf{Sets and indices.}

\begin{itemize}
\tightlist
\item
  \(m \in \mathcal{M}\): machines.
\item
  \(b \in \mathcal{B}\): blocks.
\item
  \(s=(d,\sigma) \in \mathcal{S}\): shift slots indexed by day
  \(d \in \mathcal{D}\) and shift label \(\sigma\), ordered by day and,
  within a day, by the scenario's shift order;
  \(\operatorname{prev}(s)\) is the preceding slot (the previous shift
  of the same day when there is one) and \(s'\prec s\) means slot \(s'\)
  precedes \(s\).
\item
  \(\mathcal{B}^{\text{seq}} \subseteq \mathcal{B}\): blocks with an
  explicit harvest system (\texttt{harvest\_system\_id}); blocks in
  \(\mathcal{B}\setminus\mathcal{B}^{\text{seq}}\) carry no sequencing
  obligations.
\item
  \(\mathcal{R}_b\): ordered machine roles required by the harvest
  system assigned to block \(b\in\mathcal{B}^{\text{seq}}\)
  (\(\mathcal{R}_b=\emptyset\) otherwise).
\item
  \(\mathcal{L}\): landings.
\item
  \(\mathcal{P}^{\text{inv}} \subseteq \{(r,b): r \in \mathcal{R}_b\}\):
  role-block pairs with upstream prerequisites.
\item
  \(\mathcal{P}^{\text{stg}} = \{(u,b): u \in \mathcal{U}_{r,b} \text{ for some } (r,b)\in\mathcal{P}^{\text{inv}}\}\):
  upstream role-block pairs whose output is staged for downstream roles;
  \(\mathcal{N}_{u,b} = \{r: u\in\mathcal{U}_{r,b}\}\) are the
  downstream roles of \(u\) on block \(b\).
\item
  \(\mathcal{P}^{\text{hs}} \subseteq \mathcal{P}^{\text{inv}}\): pairs
  with a head start of \(\beta_{r,b}>0\) shifts
  (\texttt{role\_headstart\_shifts}).
\item
  \(\mathcal{P}^{\text{load}} \subseteq \{(r,b): r \in \mathcal{R}_b\}\):
  loader role-block pairs;
  \(\mathcal{P}^{\text{thr}} = \{(r,b)\in\mathcal{P}^{\text{load}}\cap\mathcal{P}^{\text{inv}}: q^{\text{batch}}_{b}>0,\ W_b>0\}\):
  loaders subject to the truckload threshold.
\item
  \(\mathcal{P}^{\text{act}} = \mathcal{P}^{\text{hs}} \cup \mathcal{P}^{\text{thr}}\):
  role-block pairs whose production is gated by an activation binary.
\item
  \(\mathcal{P}^{\text{cap}} \subseteq \{(r,b): r \in \mathcal{R}_b\setminus\mathcal{T}_b\}\):
  non-terminal pairs with a per-slot remaining-volume cap; it contains
  every non-terminal pair of systems with a role feeding several roles
  or with several terminal roles, and, in the other systems, the pairs
  whose carried-in state violates
  \(R_{r,b} + \sum_{v\in\pi_{r,b}} \bar{I}_{v,b} \le W_b\)
  (\(\pi_{r,b}\): \(r\) and the non-terminal roles on its path to the
  terminal role). The cap is implied by the other constraints for all
  remaining pairs, and \(\mathcal{P}^{\text{cap}}=\emptyset\) without an
  initial state in linear and joining systems.
\item
  \(\mathcal{S}^{\text{tail}}_b \subseteq \mathcal{S}\), \(b\) with a
  loader in \(\mathcal{P}^{\text{thr}}\) and
  \(W_b > q^{\text{batch}}_b\): slots in which the remaining block
  volume can have fallen to one truckload,
  i.e.~\(\bar{D}_{b,s} > W_b - q^{\text{batch}}_b\), where
  \(\bar{D}_{b,s} = \min\{W_b, \sum_{s'\prec s}\sum_{t\in\mathcal{T}_b}\sum_{m\in\mathcal{M}(t)} A_{m,s'}\mathbf{1}^{\text{window}}_{b,d(s')}\bar{p}_{mb}\}\)
  bounds the terminal output delivered before \(s\).
\item
  \(s_1 \in \mathcal{S}\): first shift slot of the horizon.
\item
  \(\mathcal{M}^{0} \subseteq \mathcal{M}\) (\textbf{INIT}): machines
  with a known initial block \(b^{0}_m\) (optional initial state; empty
  by default); machines in \(\mathcal{M}\setminus\mathcal{M}^{0}\) start
  \emph{unplaced}.
\item
  \(\mathcal{K}\) (\textbf{E13}): locked assignments
  \(k=(m_k,b_k,d_k,\sigma_k)\), where \(\sigma_k\) is a shift label or
  empty (whole day);
  \(\mathcal{S}_k = \{s=(d_k,\sigma) \in \mathcal{S} : \sigma_k \text{ empty or } \sigma=\sigma_k\}\)
  (empty by default).
  \(\mathcal{K}_{b,s} = \{m_k: k\in\mathcal{K},\ b_k=b,\ s\in\mathcal{S}_k,\ \chi_k A_{m_k,s}=1\}\)
  are the machines locked to block \(b\) in slot \(s\).
\end{itemize}

\textbf{Parameters.}

\begin{itemize}
\tightlist
\item
  \(\bar{p}_{mb}\): production rate for machine \(m\) on block \(b\)
  (units per shift).
\item
  \(W_b\): required total block production volume.
\item
  \(A_{m,s} \in \{0,1\}\): machine availability for shift \(s\):
  \(A_{m,s}=0\) when the machine's day or shift calendar marks it
  unavailable or when \(s\) falls in a timeline blackout window, 1
  otherwise. Blackouts are fleet-wide: on every day of a blackout window
  every machine is unavailable in every slot of that day (all slots of
  \(\mathcal{S}\) on the day and the machine's own shift-calendar
  shifts), also when only some machines have shift-calendar entries.
\item
  \(\mathbf{1}^{\text{window}}_{b,d} \in \{0,1\}\): block window
  indicator (1 if day \(d\) is within block \(b\) window).
\item
  \(\omega^{\text{prod}},\omega^{\text{mob}},\omega^{\text{trans}},\omega^{\text{land}}\):
  objective weights.
\item
  \(\delta_{m,b',b} \ge 0\): mobilization cost when machine \(m\) moves
  from block \(b'\) to block \(b \ne b'\);
  \(c_{m,b',b} = \omega^{\text{mob}}\delta_{m,b',b} + \omega^{\text{trans}}\)
  is the weighted cost of the move.
\item
  \(C_{\ell}\): assignment capacity of landing \(\ell\), the number of
  machines that may work its blocks concurrently in one shift slot
  (\texttt{Landing.daily\_capacity});
  \(K_{\ell} = |\mathcal{M}| - C_{\ell}\) bounds the machines beyond
  capacity in a slot.
\item
  \(\ell(b)\): landing associated with block \(b\).
\item
  \(\mathcal{U}_{r,b}\): upstream roles that must feed role \(r\) on
  block \(b\).
\item
  \(B_{r,b}\): head-start buffer volume that every upstream role of
  \(r\) must have staged before \(r\) may produce on block \(b\):
  \(B_{r,b}=\beta_{r,b}\sum_{u\in\mathcal{U}_{r,b}}\sum_{m\in\mathcal{M}(u)}\bar{p}_{mb}\)
  for a head start of \(\beta_{r,b}\) shifts
  (\texttt{role\_headstart\_shifts}, 0 by default). When no machine of
  an upstream role has a positive rate on \(b\), the role's own capacity
  is used instead, \(B_{r,b}=\beta_{r,b}\,Q_{r,b}\). For a join the
  buffer, computed from the summed rates of all upstream roles, is
  required of each upstream role (\textbf{E8b}), so a slower upstream
  role needs more than \(\beta_{r,b}\) shifts to stage it.
\item
  \(Q_{r,b}=\sum_{m\in\mathcal{M}(r)}\bar{p}_{mb}\): role production
  capacity per shift (1 when the role has no machine with a positive
  rate on \(b\); used for the activation linearization).
\item
  \(q^{\text{batch}}_{b}\): truckload (loader batch) volume of block
  \(b\)'s harvest system (\texttt{loader\_batch\_volume\_m3}, 30 m³ by
  default).
\item
  \(\mathcal{T}_b \subseteq \mathcal{R}_b\): terminal roles for block
  \(b\) (roles credited in block completion objective terms).
\item
  \(\bar{I}_{u,b} \ge 0\) (\textbf{INIT}): initial staged volume output
  by role \(u\) on block \(b\) and not yet consumed downstream (optional
  initial state \texttt{staged\_inventory}; 0 by default).
\item
  \(R^{0}_{r,b} \ge 0\) (\textbf{INIT}): carried-in remaining volume
  role \(r\) may still output on block \(b\) (optional initial state
  \texttt{role\_remaining}; \(W_b\) by default).
\item
  \(R_{r,b} = \min(R^{0}_{r,b}, W_b)\): remaining volume role \(r\) may
  output on block \(b\) (every role handles the same wood, so no role
  can output more than the block still holds).
\item
  \(b^{0}_m\) (\textbf{INIT}): block machine \(m\in\mathcal{M}^{0}\)
  occupied in its last worked slot before the horizon (optional initial
  state \texttt{last\_block\_id}).
\end{itemize}

\textbf{Decision variables.}

\begin{itemize}
\tightlist
\item
  \(x_{m,b,s} \in \{0,1\}\): 1 if machine \(m\) is assigned to block
  \(b\) in shift \(s\).
\item
  \(p_{m,b,s} \ge 0\): production by machine \(m\) on block \(b\) in
  shift \(s\).
\item
  \(z_{r,b,s} \ge 0\): aggregated role-level production for role \(r\)
  on block \(b\) in shift \(s\).
\item
  \(y_{m,b',b,s} \ge 0\), \(b' \ne b\): machine \(m\) moves in slot
  \(s\) from its \emph{position} \(b'\) (the block of its last worked
  slot before \(s\), or \(b^{0}_m\)) to block \(b\), which it works in
  \(s\).
\item
  \(\eta_{m,b,s} \ge 0\): machine \(m\) keeps position \(b\) through
  slot \(s\) (it is idle or works \(b\) again).
\item
  \(\phi_{m,b,s} \ge 0\), \(m\notin\mathcal{M}^{0}\): the first worked
  slot of machine \(m\) is \(s\), on block \(b\); \(\nu_{m,s} \ge 0\),
  \(m\notin\mathcal{M}^{0}\): machine \(m\) has not worked up to and
  including \(s\) (both \(\equiv 0\) for \(m\in\mathcal{M}^{0}\)).
\item
  \(\pi_{m,b,s} = \eta_{m,b,s} + \sum_{b'\ne b} y_{m,b',b,s} + \phi_{m,b,s}\):
  machine \(m\) holds position \(b\) after slot \(s\) (notation);
  \(\pi_{m,b,\operatorname{prev}(s_1)} := \mathbf{1}[m\in\mathcal{M}^{0},\, b=b^{0}_m]\)
  and
  \(\nu_{m,\operatorname{prev}(s_1)} := \mathbf{1}[m\notin\mathcal{M}^{0}]\).
\item
  \(I^{\text{start}}_{u,b,s} \ge 0\),
  \((u,b)\in\mathcal{P}^{\text{stg}}\): output of upstream role \(u\) on
  block \(b\) staged for its downstream roles at the start of shift
  \(s\).
\item
  \(I_{u,b,s} \ge 0\): the same staged volume at the end of shift \(s\).
\item
  \(g_{r,b,s} \in \{0,1\}\), \((r,b)\in\mathcal{P}^{\text{act}}\): role
  activation indicator; it gates production only.
\item
  \(h_{r,b,s} \in \{0,1\}\), \((r,b)\in\mathcal{P}^{\text{hs}}\): 1 only
  if every upstream role of \(r\) on block \(b\) has output its whole
  carried-in remaining volume before slot \(s\) (the buffer can no
  longer grow and is waived).
\item
  \(\lambda_{b,s} \in \{0,1\}\), \(s\in\mathcal{S}^{\text{tail}}_b\): 1
  only once the remaining volume of block \(b\) is at most one truckload
  (selects the active term of the loader threshold).
\item
  \(D_{b,s} = \sum_{t\in\mathcal{T}_b}\sum_{s'\preceq s} z_{t,b,s'}\):
  terminal output delivered on block \(b\) up to and including slot
  \(s\) (notation for a cumulative sum;
  \(D_{b,\operatorname{prev}(s_1)} := 0\)).
\item
  \(L_b \ge 0\): leftover unmet block volume slack.
\item
  \(S_{\ell,s,k} \in [0,1]\), \(k = 1,\dots,K_{\ell}\): unit landing
  surplus slack for the \(k\)-th machine beyond capacity on landing
  \(\ell\) in slot \(s\) (only when \(\omega^{\text{land}} > 0\)).
\end{itemize}

\textbf{Objective (OBJ).}

FHOPS maximizes weighted terminal production and subtracts penalty
terms:

\[
\begin{aligned}
\max\; &\omega^{\text{prod}}\!\sum_{b\in\mathcal{B}}\sum_{r\in\mathcal{T}_b}\sum_{s\in\mathcal{S}} z_{r,b,s}
- \omega^{\text{prod}}\!\sum_{b\in\mathcal{B}} L_b \\
&- \omega^{\text{land}}\!\sum_{\ell\in\mathcal{L}}\sum_{s\in\mathcal{S}}\sum_{k=1}^{K_{\ell}} k\, S_{\ell,s,k} \\
&- \sum_{m\in\mathcal{M}}\sum_{s\in\mathcal{S}}\sum_{b'\ne b}
\left(\omega^{\text{mob}}\,\delta_{m,b',b} + \omega^{\text{trans}}\right) y_{m,b',b,s}.
\end{aligned}
\]

The last line charges every move of a machine (\textbf{E6}): working a
block other than its position, including the move from \(b^{0}_m\) into
the machine's first worked slot. Staying on a block costs nothing, and a
machine without initial block (\(m\notin\mathcal{M}^{0}\)) moves for
free into its first worked slot.

For blocks without terminal roles (\(\mathcal{T}_b=\emptyset\), in
particular blocks outside \(\mathcal{B}^{\text{seq}}\)) the production
reward and the block balance use the machine-level sum
\(\sum_{m}\sum_{s} p_{m,b,s}\) in place of
\(\sum_{r\in\mathcal{T}_b}\sum_s z_{r,b,s}\).

\textbf{Earliness tie-break (OBJ2; optional, default in rolling-horizon
windows).} OBJ does not depend on when work is done inside the horizon,
so plans that shift production between slots tie. In a rolling-horizon
window, a tied optimum may defer work past the lock span and lock idle
days. With the earliness option, a second stage maximizes the
production-weighted earliness

\[
E=\sum_{s\in\mathcal{S}} w_s \sum_{m\in\mathcal{M}}\sum_{b\in\mathcal{B}} p_{m,b,s},
\qquad w_s=\frac{|\mathcal{S}|-k_s}{|\mathcal{S}|},
\]

where \(k_s\in\{0,\dots,|\mathcal{S}|-1\}\) is the position of slot
\(s\) in the slot order. This stage keeps every constraint of
\textbf{E1}--\textbf{E13} and adds \(\text{OBJ}\ge z_1-\tau\), where
\(z_1\) is the OBJ value of the stage-1 solution and
\(\tau=10^{-6}\max(1,|z_1|)\). Stage 2 is warm-started from the stage-1
solution, so its returned plan has \(\text{OBJ}\ge z_1-\tau\). If stage
1 is optimal, the returned plan is therefore optimal for OBJ within
\(\tau\), which is 100 times tighter than HiGHS's default relative MIP
gap (\(10^{-4}\)). The reported objective is OBJ; \(E\) is reported
separately. A single weighted objective \(\text{OBJ}+\varepsilon E\) is
not used, for two reasons. First, production is continuous and the data
are arbitrary reals, so no data-independent \(\varepsilon>0\) is
guaranteed to stay below the smallest positive OBJ difference between
plans: for any \(\varepsilon\) there are data for which the weighted
optimum is not OBJ-optimal. Second, an \(\varepsilon\) small enough to
be harmless in practice (\(\varepsilon E\ll 10^{-4}|z_1|\)) is below the
solver's relative gap, so the solver would stop before it acts on the
tie-break. The cost is a second solve. Rolling-horizon MILP windows
enable the option by default, except windows whose lock span covers the
whole window. Standalone solves (\texttt{fhops\ solve-mip-operational},
\texttt{solve\_operational\_milp}) do not enable it by default, because
the published single-horizon optimum and its solve time stay unchanged;
\texttt{-\/-earliness} or \texttt{earliness=True} turns it on.

\textbf{Constraints.}

Machine assignment feasibility, including calendar availability and
timeline blackouts through \(A_{m,s}\) (\textbf{E1}):

\[
\sum_{b\in\mathcal{B}} x_{m,b,s} \le A_{m,s}
\qquad \forall m\in\mathcal{M},\; s\in\mathcal{S}.
\]

Role compatibility, machines can only work roles allowed by the block's
assigned harvest system (\textbf{E2}):

\[
x_{m,b,s}=0 \quad \text{if } b\in\mathcal{B}^{\text{seq}} \text{ and role}(m)\notin\mathcal{R}_b.
\]

Production upper bound per assignment (\textbf{E3}):

\[
p_{m,b,s} \le \bar{p}_{mb}\,x_{m,b,s}
\qquad \forall m,b,s.
\]

Block window enforcement (\textbf{E4}):

\[
x_{m,b,s}=0 \quad \text{if } \mathbf{1}^{\text{window}}_{b,d}=0 \text{ for } s=(d,\sigma).
\]

Role-production aggregation (\textbf{E5}):

\[
z_{r,b,s} = \sum_{m\in\mathcal{M}(r)} p_{m,b,s}
\qquad \forall (r,b), s.
\]

Machine positions and moves (\textbf{E6}; each machine's position is a
unit flow through one layer per slot; idle slots keep the position):

\[
\pi_{m,b,\operatorname{prev}(s)} = \eta_{m,b,s} + \sum_{b''\ne b} y_{m,b,b'',s}
\qquad \forall m\in\mathcal{M},\; b\in\mathcal{B},\; s\in\mathcal{S},
\]

\[
\nu_{m,\operatorname{prev}(s)} = \nu_{m,s} + \sum_{b\in\mathcal{B}} \phi_{m,b,s}
\qquad \forall m\in\mathcal{M}\setminus\mathcal{M}^{0},\; s\in\mathcal{S},
\]

\[
\sum_{b'\ne b} y_{m,b',b,s} + \phi_{m,b,s} \;\le\; x_{m,b,s} \;\le\; \pi_{m,b,s}
\qquad \forall m\in\mathcal{M},\; b\in\mathcal{B},\; s\in\mathcal{S}.
\]

A position changes only into a worked block, and working block \(b\)
puts the whole unit of flow at \(b\). For binary \(x\), every path of a
decomposition of the flow therefore follows the machine's true position
sequence, so \(y_{m,b',b,s}=1\) exactly when machine \(m\) works \(b\)
in slot \(s\) and its last worked block before \(s\) (or \(b^{0}_m\)) is
\(b'\ne b\), also when idle slots lie in between; all other \(y\) are 0.
The move term of OBJ is thus exact and equals the mobilization and
transition accounting of the heuristics and of the playback KPIs. The
implementation builds the network only for machines that can incur a
positive move cost and only for the slots in which the machine can work
(other slots keep its position), starting arcs only at positions
reachable before the slot; when all move costs \(c_{m,b',b}\) of a
machine are equal, the arcs \(y_{m,b',b,s}\) are replaced by arcs to and
from one hub node per slot, an equivalent network with
\(O(|\mathcal{B}|)\) instead of \(O(|\mathcal{B}|^2)\) arcs.

Staged inventory start and balance per upstream role (\textbf{E7}; each
downstream role consumes its own output from the staged output of
\textbf{every} upstream role, so a role with several upstream roles can
only process what each of them has staged):

\[
I^{\text{start}}_{u,b,s}=
\begin{cases}
\bar{I}_{u,b}, & s = s_1\\
I_{u,b,\operatorname{prev}(s)}, & \text{otherwise}
\end{cases}
\qquad \forall (u,b)\in\mathcal{P}^{\text{stg}}, s,
\]

\[
I_{u,b,s}=I^{\text{start}}_{u,b,s}+z_{u,b,s}-\sum_{r\in\mathcal{N}_{u,b}} z_{r,b,s}
\qquad \forall (u,b)\in\mathcal{P}^{\text{stg}}, s,
\]

\[
\sum_{r\in\mathcal{N}_{u,b}} z_{r,b,s} \le I^{\text{start}}_{u,b,s}
\qquad \forall (u,b)\in\mathcal{P}^{\text{stg}}, s.
\]

Staged output is therefore available downstream from the next shift
slot. For a linear chain (\(|\mathcal{U}_{r,b}|=|\mathcal{N}_{u,b}|=1\))
these are the FHOPS 1.0.0 inventory equations of the downstream role.
The downstream roles of a fork (\(|\mathcal{N}_{u,b}|>1\)) split the
staged output of \(u\): each unit is consumed by one of them. A join
(\(|\mathcal{U}_{r,b}|>1\)) consumes each unit of its output from the
pool of every upstream role. Hence, without carried-in staged volume, a
fork that joins again (a diamond \(u\to\{r_1,r_2\}\to t\)) delivers at
most half of the output of \(u\), i.e.~at most \(W_b/2\).

Activation, production gating (\textbf{E8a}):

\[
\begin{aligned}
&z_{r,b,s} \le Q_{r,b}\,g_{r,b,s},
\qquad
g_{r,b,s} \le \sum_{m\in\mathcal{M}(r)} x_{m,b,s},\\
&\sum_{m\in\mathcal{M}(r)\setminus\mathcal{K}_{b,s}} x_{m,b,s} \le |\mathcal{M}(r)\setminus\mathcal{K}_{b,s}|\, g_{r,b,s}
\qquad \forall (r,b)\in\mathcal{P}^{\text{act}}, s.
\end{aligned}
\]

An assigned unlocked machine activates its role; a machine locked to the
block may stay idle (\(x=1\), \(p=0\)) without activating it, so a lock
never forces production that the staged volumes cannot support.

Head-start buffer (\textbf{E8b}; with
\(I_{u,b,\operatorname{prev}(s_1)} := \bar{I}_{u,b}\)):

\[
I_{u,b,\operatorname{prev}(s)} \ge B_{r,b}\,\left(g_{r,b,s}-h_{r,b,s}\right)
\qquad \forall (r,b)\in\mathcal{P}^{\text{hs}},\; u\in\mathcal{U}_{r,b},\; s.
\]

Head-start waiver (\textbf{E8c}; the buffer is waived once every
upstream role has output its carried-in remaining volume, so blocks
smaller than a buffer can still be finished):

\[
\sum_{s'\prec s} z_{u,b,s'} \ge R^{0}_{u,b}\,h_{r,b,s}
\qquad \forall (r,b)\in\mathcal{P}^{\text{hs}},\; u\in\mathcal{U}_{r,b},\; s.
\]

Loader truckload threshold (\textbf{E9}): a producing loader needs one
truckload, or the whole remaining block volume when it is smaller,
staged by every upstream role at the start of the slot,

\[
I_{u,b,\operatorname{prev}(s)} \ge \min\!\left(q^{\text{batch}}_{b},\; W_b - D_{b,\operatorname{prev}(s)}\right) g_{r,b,s}
\qquad \forall (r,b)\in\mathcal{P}^{\text{thr}},\; u\in\mathcal{U}_{r,b},\; s,
\]

linearized exactly as follows (all coefficients are data;
\(W_b - D_{b,\operatorname{prev}(s)} \in [0, W_b]\)):

\[
\begin{aligned}
&\text{if } W_b > q^{\text{batch}}_b \text{ and } s\notin\mathcal{S}^{\text{tail}}_b:\\
&\qquad I_{u,b,\operatorname{prev}(s)} \ge q^{\text{batch}}_{b}\, g_{r,b,s},\\
&\text{if } W_b \le q^{\text{batch}}_b:\\
&\qquad I_{u,b,\operatorname{prev}(s)} + D_{b,\operatorname{prev}(s)} \ge W_b\, g_{r,b,s},\\
&\text{if } s\in\mathcal{S}^{\text{tail}}_b:\\
&\qquad I_{u,b,\operatorname{prev}(s)} \ge q^{\text{batch}}_{b}\,(g_{r,b,s}-\lambda_{b,s}),\\
&\qquad I_{u,b,\operatorname{prev}(s)} + D_{b,\operatorname{prev}(s)} \ge q^{\text{batch}}_{b}\, g_{r,b,s} + (W_b - q^{\text{batch}}_{b})\,\lambda_{b,s},
\end{aligned}
\]

\[
\begin{aligned}
&D_{b,\operatorname{prev}(s)} \ge (W_b - q^{\text{batch}}_{b})\,\lambda_{b,s}
&& \forall s\in\mathcal{S}^{\text{tail}}_b,\\
&\lambda_{b,s} \ge \lambda_{b,\operatorname{prev}(s)}
&& \forall s \text{ with } s,\operatorname{prev}(s)\in\mathcal{S}^{\text{tail}}_b.
\end{aligned}
\]

\(\lambda_{b,s}=1\) is only possible once at most one truckload remains;
it then relaxes the threshold to the remaining volume. Outside
\(\mathcal{S}^{\text{tail}}_b\) the remaining volume provably exceeds a
truckload, so no binary is needed there.

Block completion balance with leftover slack (\textbf{E10}):

\[
\sum_{r\in\mathcal{T}_b}\sum_{s\in\mathcal{S}} z_{r,b,s} + L_b = W_b
\qquad \forall b\in\mathcal{B}.
\]

Landing capacity per shift slot, the machines working a landing's blocks
concurrently (\textbf{E11}):

\[
\sum_{b:\,\ell(b)=\ell}\sum_{m\in\mathcal{M}} x_{m,b,s}
\le
\begin{cases}
\max\{C_{\ell},\, N^{\text{lock}}_{\ell,s}\} & \text{if } \omega^{\text{land}} = 0,\\
C_{\ell} + \sum_{k=1}^{K_{\ell}} S_{\ell,s,k} & \text{if } \omega^{\text{land}} > 0,
\end{cases}
\qquad \forall \ell\in\mathcal{L},\; s\in\mathcal{S},
\]

where
\(N^{\text{lock}}_{\ell,s} = \sum_{b:\,\ell(b)=\ell}|\mathcal{K}_{b,s}|\)
is the number of machines locked to the landing's blocks in slot \(s\).
With \(\omega^{\text{land}}=0\) the capacity is hard; a slot in which
locks alone exceed it keeps the locked machines, admits no other
machine, and is reported as a warning (the heuristics charge the same
unavoidable overload). With \(\omega^{\text{land}}>0\) the marginal
price \(k\,\omega^{\text{land}}\) of the slack pieces increases, so an
optimal solution uses \(S_{\ell,s,1},\dots,S_{\ell,s,e}\) for \(e\)
machines beyond capacity and pays \(\omega^{\text{land}}\,e(e+1)/2\):
the \(k\)-th machine beyond capacity in a slot costs
\(k\,\omega^{\text{land}}\), as in the heuristics' evaluation.

Remaining role output, no role can handle more wood than the block still
holds (\textbf{E12}):

\[
\begin{aligned}
&\sum_{s\in\mathcal{S}} z_{r,b,s} \le R_{r,b}
&& \forall b\in\mathcal{B}^{\text{seq}},\; r\in\mathcal{R}_b,\\
&z_{r,b,s} + D_{b,s} \le W_b
&& \forall (r,b)\in\mathcal{P}^{\text{cap}},\; s.
\end{aligned}
\]

Locked assignments (\textbf{E13}; a lock without a shift label pins
every available shift of its day; a lock with a shift label pins only
that slot):

\[
x_{m_k,b,s} = \chi_k\,A_{m_k,s}\,\mathbf{1}[b=b_k]
\qquad \forall k\in\mathcal{K},\; s\in\mathcal{S}_k,\; b\in\mathcal{B},
\]

where \(\chi_k=0\) when the lock contradicts the model (day outside the
block window, or a machine whose role is not in the block's harvest
system) and \(\chi_k=1\) otherwise. Such locks are rejected by scenario
validation; a lock that still reaches the model pins the machine to idle
and is reported as a warning instead of making the model infeasible. A
second lock on an already locked machine slot is ignored with a warning.

Domain restrictions (\textbf{D1}):

\[
x, g, h, \lambda \in \{0,1\},\quad p,z,y,\eta,\phi,\nu,I^{\text{start}},I,L \ge 0,\quad S \in [0,1].
\]

\textbf{Initial state (INIT).} The optional
\texttt{Scenario.initial\_state} supplies \(\bar{I}_{u,b}\)
(\textbf{E7}, \textbf{E8b}), \(R^{0}_{r,b}\) (\textbf{E8c},
\textbf{E12}), and \(b^{0}_m\) with \(\mathcal{M}^{0}\) (\textbf{E6});
the rolling-horizon driver uses it to carry state from one planning
window to the next. Without \texttt{Scenario.initial\_state} and
\texttt{Scenario.locked\_assignments} (\(\bar{I}\equiv 0\);
\(R^{0}_{r,b}= R_{r,b}= W_b\);
\(\mathcal{M}^{0}=\mathcal{K}=\emptyset\)) the initial-state and lock
terms vanish.

\textbf{Changes from FHOPS v1.0.0 (1.0.1).} The following corrections
make every MILP plan physically feasible and replayable by the playback
sequencing tracker without violations: (i) (\textbf{E12}) the
remaining-output cap applies to every role with
\(R_{r,b}=\min(R^{0}_{r,b},W_b)\) (\(W_b\) by default), complemented by
the per-slot cap \(z_{r,b,s}+D_{b,s}\le W_b\) where the flow balances do
not imply it (v1.0.0 let upstream roles output more volume than the
block holds and used that volume to meet head-start and loader
thresholds); (ii) (\textbf{E8c}) the head-start constraint is waived
through \(h_{r,b,s}\) once every upstream role has output its carried-in
remaining volume, so blocks smaller than a buffer can still be finished;
(iii) (\textbf{E9}) the loader threshold is
\(\min(q^{\text{batch}}_b, W_b - D_{b,\operatorname{prev}(s)})\), the
volume still to deliver at the start of the slot (v1.0.0 always required
a full truckload), so the last partial truckload of a block can be
loaded; (iv) (\textbf{E7}) staged inventories are kept per upstream
role, \(I_{u,b,s}\) (v1.0.0 kept one inventory per downstream role fed
by the sum of its upstream roles' outputs, so a role with several
upstream roles could process wood that only one of them had handled);
(v) (\textbf{E8a}, \textbf{E13}) the activation binary \(g_{r,b,s}\)
gates production only and machines locked to a block may stay idle, so
locks (new in 1.0.1) never make the model infeasible; contradictory
locks are pinned to idle and reported; (vi) (\textbf{E2}) blocks without
a harvest system (\(\mathcal{B}\setminus\mathcal{B}^{\text{seq}}\)) have
no role obligations, as documented in the data contract (v1.0.0 applied
the registry's default system to them); (vii) (\textbf{E1}) timeline
blackouts, which v1.0.0 enforced only in the heuristics, set
\(A_{m,s}=0\) in every slot of a blackout day for every machine; (viii)
(\textbf{E11}) landing capacity counts the machines on a landing per
shift slot and is hard when \(\omega^{\text{land}}=0\), as in the
heuristics (v1.0.0 counted machine-shifts per day,
\(\sum_{\sigma} x_{m,b,(d,\sigma)} \le C_{\ell} + S_{\ell,d}\), with a
slack that was free at the default \(\omega^{\text{land}}=0\), so the
capacity did not bind; on single-shift scenarios the per-slot and
per-day counts coincide); (ix) (\textbf{E6}) a move is charged when a
machine works a block other than its position, its last worked block or
\(b^{0}_m\), through the position network above: staying on a block
costs nothing (v1.0.0 also charged \(\omega^{\text{trans}}\) for
consecutive slots on the same block), and a move across idle slots, or
from \(b^{0}_m\) into a first worked slot after \(s_1\), is charged
(v1.0.0 linked only consecutive slots through binaries
\(y_{m,b',b,s}\ge x_{m,b',\operatorname{prev}(s)}+x_{m,b,s}-1\), so
idling one slot avoided the move cost that the heuristics and the KPIs
charge); (x) (old \textbf{E9}) the loader batching variables of v1.0.0
(\(z_{r,b,s}=q^{\text{batch}}_{b}n_{r,b,s}+u_{r,b,s}\) with
\(n_{r,b,s}\in\mathbb{Z}_{\ge0}\),
\(0\le u_{r,b,s}\le q^{\text{batch}}_{b}\)) were removed because every
\(z_{r,b,s}\ge 0\) satisfies them; the truckload rule is the loader
threshold above. Playback, the heuristics, and the rolling-horizon
carry-forward apply the same rules: staged output is available from the
next shift slot, buffers and truckload thresholds are staged volume at
the start of the slot, and production is capped by \(R_{r,b}\) and by
the volume the block still holds. For linear harvest systems without
loaders, without locks, and with no (or a consistent) initial state,
(iii)--(v) leave the equations unchanged apart from indexing the staged
inventory by the upstream instead of the downstream role. The initial
state (\textbf{INIT}), locked assignments (\textbf{E13}), and the
earliness stage (\textbf{OBJ2}) are new in 1.0.1 and vanish by default.

\textbf{Implementation mapping (equation blocks to code).}

Primary modules are \texttt{operational.py}, \texttt{data.py}, and
\texttt{driver.py} under \texttt{src/fhops/model/milp}. Regression tests
live under \texttt{tests/model} (builder, moves, landing capacity,
blackouts, earliness, driver) and cover the MILP builder plus driver
integration.

\begin{longtable}[]{@{}
  >{\raggedright\arraybackslash}p{(\linewidth - 2\tabcolsep) * \real{0.1600}}
  >{\raggedright\arraybackslash}p{(\linewidth - 2\tabcolsep) * \real{0.8400}}@{}}
\toprule\noalign{}
\begin{minipage}[b]{\linewidth}\raggedright
Block
\end{minipage} & \begin{minipage}[b]{\linewidth}\raggedright
Pyomo objects and rules
\end{minipage} \\
\midrule\noalign{}
\endhead
\bottomrule\noalign{}
\endlastfoot
\textbf{OBJ} & \texttt{model.objective} and the objective terms around
\texttt{prod\_weight}, \texttt{landing\_weight},
\texttt{mobilisation\_weight}, \texttt{transition\_weight} (move costs
\(c_{m,b',b}\); \(b^{0}_m\) from
\texttt{bundle.initial\_machine\_block}) \\
\textbf{OBJ2} & \texttt{earliness\_expression(...)} in
\texttt{fhops.model.milp.operational} (\(E\)); second stage
\texttt{\_solve\_earliness\_stage(...)} in
\texttt{fhops.model.milp.driver} (\texttt{model.earliness\_floor},
\texttt{model.earliness\_objective};
\texttt{solve\_operational\_milp(...,\ earliness=True)}, result key
\texttt{earliness}) \\
\textbf{INIT} & \texttt{build\_operational\_bundle(...)} in
\texttt{fhops.model.milp.data} flattens \texttt{Scenario.initial\_state}
and \texttt{Scenario.locked\_assignments} into the bundle
(\texttt{bundle.initial\_staged\_inventory},
\texttt{bundle.initial\_role\_remaining},
\texttt{bundle.initial\_machine\_block},
\texttt{bundle.locked\_assignments}) \\
\textbf{E1} & \texttt{model.machine\_capacity}
(\texttt{machine\_capacity\_rule}; \(A_{m,s}\) from the calendars and
\texttt{bundle.blackout\_slots}, built by
\texttt{build\_blackout\_slots(...)} and shared with the heuristics) \\
\textbf{E2} & \texttt{model.role\_compatibility}
(\texttt{role\_compatibility\_rule}; skipped for
\texttt{bundle.unsequenced\_blocks}) \\
\textbf{E3} & \texttt{model.production\_cap}
(\texttt{prod\_cap\_rule}) \\
\textbf{E4} & \texttt{model.block\_windows} (\texttt{window\_rule}) \\
\textbf{E5} & \texttt{model.role\_prod\_balance}
(\texttt{role\_prod\_balance\_rule}) \\
\textbf{E6} & \texttt{model.position\_balance} (\(\eta\) =
\texttt{model.stay}, \(y\) = \texttt{model.y} or the hub arcs
\texttt{model.depart}/\texttt{model.arrive} with
\texttt{model.hub\_balance} for machines with equal move costs),
\texttt{model.unplaced\_balance} (\(\phi\) = \texttt{model.first},
\(\nu\) = \texttt{model.unplaced}), \texttt{model.move\_requires\_work},
\texttt{model.work\_sets\_position} \\
\textbf{E7} & \texttt{model.inventory\_start\_eq} (indexed by upstream
role-block pairs \texttt{model.InventoryPairs}; first slot uses
\(\bar{I}_{u,b}\) from \texttt{bundle.initial\_staged\_inventory}),
\texttt{model.inventory\_balance}, \texttt{model.inventory\_guard} \\
\textbf{E8a} & \texttt{model.activation\_prod},
\texttt{model.role\_active\_upper} (locked machines excluded),
\texttt{model.role\_active\_lower} \\
\textbf{E8b} & \texttt{model.head\_start} (one row per upstream role);
\(B_{r,b}\) from \texttt{headstart\_buffer\_volumes(...)} in
\texttt{fhops.model.milp.data} \\
\textbf{E8c} & \texttt{model.upstream\_done\_link} (cumulative output
\texttt{model.role\_cumulative\_eq}) \\
\textbf{E9} & \texttt{model.loader\_threshold},
\texttt{model.loader\_threshold\_tail},
\texttt{model.loader\_tail\_reached},
\texttt{model.loader\_tail\_monotone} (\(\lambda\) =
\texttt{model.loader\_tail}; \(D_{b,s}\) from
\texttt{model.role\_cumulative} of the terminal roles) \\
\textbf{E10} & \texttt{model.block\_balance}
(\texttt{block\_balance\_rule}) + \texttt{model.leftover} \\
\textbf{E11} & \texttt{model.landing\_capacity}
(\texttt{landing\_capacity\_rule}, indexed by landing and slot; locked
overloads from \texttt{\_landing\_slot\_capacities}) +
\texttt{model.landing\_surplus} (unit pieces
\texttt{model.LandingSurplusIndex}, only when
\(\omega^{\text{land}}>0\)) \\
\textbf{E12} & \texttt{model.role\_remaining\_cap} (from
\texttt{bundle.initial\_role\_remaining}, capped at \(W_b\)), per-slot
cap \texttt{model.role\_slot\_remaining} \\
\textbf{E13} & \texttt{model.locked\_assignment} (from
\texttt{bundle.locked\_assignments}, resolved by
\texttt{resolve\_locked\_slots(...)} in \texttt{fhops.model.milp.data};
warnings in the solve result) \\
\textbf{D1} & Domains of \texttt{model.x}, \texttt{model.prod},
\texttt{model.role\_prod}, \texttt{model.y}, \texttt{model.stay},
\texttt{model.depart}, \texttt{model.arrive}, \texttt{model.first},
\texttt{model.unplaced}, \texttt{model.inventory\_start},
\texttt{model.inventory}, \texttt{model.role\_cumulative},
\texttt{model.role\_active}, \texttt{model.upstream\_done},
\texttt{model.loader\_tail}, \texttt{model.leftover},
\texttt{model.landing\_surplus} \\
\end{longtable}

This formulation is the canonical mathematical reference for FHOPS
operational MILP documentation and thesis-level reporting.
